\chapter{Global transport}
\label{impl:global}
\rev{This chapter mirrors Chapter~\ref{ch:global}. Built and recorded here: the two
mean-squared-displacement estimators and their audit, the filtered-variation family that
replaces them as a source of exponents, the moment spectrum, the non-Gaussian parameter, the
velocity autocorrelation and its tail, and the persistence runs. Three observables the body
catalogues are defined and not measured, and Table~\ref{tab:obs-status} says which.}

\section{Observables on the whole trajectory}
\label{impl:obs-global}
\rev{The mean-squared displacement of Sec.~\ref{sec:obs-global} is implemented and
unit-tested, as are four more of that section's observables; the shape ones are documented
below and the three that are not measured are named in Table~\ref{tab:obs-status}. The
estimators live in
\path{soaring.analysis.observables.transport} and the figure, the curve table and the
\texttt{\textbackslash StatMsd*} macros come from one pass of
\path{generate_msd_figure.py} (Table~\ref{tab:impl-repro}). A second pass,
\path{audit_msd.py}, keeps what the first one averages away, and
\path{audit_msd_report.py} reduces it into the \texttt{\textbackslash StatAudit*}
macros; both are described at the end of this section.}

\subsection{\rev{The two estimators}}
\rev{\path{MSDAccumulator} builds the ensemble average and \path{TAMSDAccumulator} the
time average. Both accumulate by streaming rather than by loading: the \path{fixes} table
runs to ${\sim}1.4\times10^{9}$ rows over the two archives, and neither estimator needs
more than one flight in memory at a time.}

\paragraph{\rev{Coverage.}} \rev{A flight answers a lag $t$ with the fix nearest to it,
and only if two conditions hold: that fix lies within $\max(\SI{0.5}{\second},\Delta
t/2)$ of $t$, and $t\ge\Delta t$, where $\Delta t$ is the median inter-fix time of the
flight's own concatenated segments. The second condition is the one that is easy to omit
and expensive to omit. Without it a \SI{5}{\second} logger asked for a lag of
\SI{1}{\second} answers with its fix at $t=0$---a displacement of exactly zero---and a
slow-cadence minority drags the short-lag MSD down until the grid crosses their step. The
same rule appears in the time-averaged estimator, which refuses any lag below its
segment's own step. A lag falling inside a gap between two retained segments is answered
by neither side, since the split bound (Sec.~\ref{sec:uniform}) makes every gap wider than
the tolerance.}

\paragraph{\rev{The lag grid.}} \rev{\path{log_lag_grid} spaces \num{90} points
geometrically from \SI{1}{\second} to \SI{43200}{\second}, the duration bound of the
flight-level filter, and snaps them to the second. The snapping is what makes the short
end usable: unsnapped, lags of \num{1.00}, \num{1.13} and \num{1.27} seconds all resolve
to the same fix at a \SI{1}{\hertz} cadence, and the curve there is a staircase of
repeated values whose local slope alternates between zero and a spurious spike. Snapping
merges them instead, which is why the grid delivers \num{80} distinct lags rather than
\num{90}, and why the lags below \SI{8}{\second} are consecutive integers. At long lags
the rounding is a part in $10^4$. The consecutive-integer run is also what lets the
coverage artefact of Sec.~\ref{sec:obs-global} have a period; that is the cost of the
snapping, it is measured in Sec.~\ref{sec:obs-global} by splitting the archive by cadence,
and it sits an order of magnitude below the lower end of every fit range.}

\paragraph{\rev{The time average.}} \rev{\path{time_averaged_msd} is computed within a
segment, never across the gap that ends it. The direct double loop is $O(n^2)$ and
hopeless at $10^4$ samples per segment; writing $D_j=|\mathbf{r}_j|^2$, the identity
$\overline{\delta^2}(k)=\langle D_j+D_{j+k}\rangle_j-2\langle
\mathbf{r}_j\cdot\mathbf{r}_{j+k}\rangle_j$ turns the second term into an
autocorrelation, which the FFT gives in $O(n\log n)$. The curve is cut at half a segment's
length, beyond which the average over starting times has too few of them to be one.}

\subsection{\rev{The fit, and what is quoted from it}}
\paragraph{\rev{The range.}} \rev{The lower end is an argument, not a consequence: at
\SI{120}{\second} it sits above the thermalling period, so what the slope measures is
transport rather than the circling superimposed on it, and above the whole region the
short-lag artefacts occupy. The upper end is a consequence.
\path{coverage_limited_range} walks the count-per-lag curve and cuts where it has fallen
to \path{min_coverage} = \num{0.25} of its value at the lower end; past that the curve is
an average over a survivorship-selected minority. Applied to the four curves this gives
the ranges quoted with each exponent.}

\paragraph{\rev{The two errors.}} \rev{\path{fit_msd_exponent} fits unweighted in log
space, which gives every decade the same say; a fit weighted by the much larger absolute
errors at long lags would be dominated by the last decade alone. It returns the
least-squares error on the slope, which is written out as
\texttt{\textbackslash StatMsd*AlphaErrOls} and is \emph{not} the quoted uncertainty.
\path{bootstrap_alpha_error} resamples \num{200} ensembles of flights with replacement,
refits each over the same range, and reports the standard deviation as
\texttt{\textbackslash StatMsd*AlphaErr}. Neither figure is attached to an exponent in the
body: both MSD exponents are quoted bare, because the bootstrap treats clustered flights as
independent records and because the range dependence is fifty times larger where both have
been measured. Sec.~\ref{sec:obs-global} makes that argument and gives the synthetic check
that establishes what each of the two errors measures; the exponents that do carry an
interval are those of Sec.~\ref{sec:transport-measure}.}

\paragraph{\rev{The cohorts.}} \rev{Three fixed-duration cohorts at $T=1$, $2$ and
\SI{4}{\hour}. Within a cohort every member is present at every lag up to $T$, so the
population is constant by construction and any remaining growth is motion; each is
therefore fitted from the range's lower end up to its own $T$, which is as far as the
guarantee reaches. The time-averaged cohorts are keyed on segment span and capped at
$T/2$, for the same reason the time-averaged curve is. A cohort is compared against the
\emph{pooled curve refitted over the cohort's own lags}, never against the pooled
exponent over the pooled range: the local slope is not constant across those two spans, so
the naive comparison charges the shape of the curve to the change of population. Both
readings are generated (\texttt{\textbackslash StatMsd*CohortSpread} and
\texttt{*CohortSpreadNaive}), and on the ensemble estimator they differ by a factor of
five, which is the size of the mistake avoided.}

\paragraph{\rev{The local slope, and why it is a regression.}}
\begin{revblock}
The panel that carries the most interpretation is the one whose estimator is easiest
to get wrong. The local slope is $\mathrm{d}\log\mathrm{MSD}/\mathrm{d}\log t$, and the
obvious way to compute it---a centred difference between neighbouring lags---divides by
the grid spacing, which on the logarithmic grid used here is about \num{0.05} decades.
The division multiplies whatever scatter the curve carries by $1/(2\times0.05)\approx4$:
a point-to-point scatter of a few tenths of a per cent in the MSD, which is about what the
ensemble's own standard error allows, becomes a slope that jitters by \num{0.05} and reads
as noise. The panel then testifies against a curve that is not in fact noisy.

\path{soaring.analysis.observables.transport.local_slope} therefore works on a window: at each lag it
takes the points within \num{0.15} decades on either side---five to six at this grid
spacing---and reports the \emph{median of the slopes of all pairs} of them, a Theil--Sen
estimator. The choice is for resistance to a single displaced lag rather than to a
heavy-tailed residual, which is the failure this curve actually has: a coverage artefact
of a few per cent at one lag tilts a least-squares window and moves a median by nothing.
The uncertainty reported with it, and drawn as a band, is the interquartile range of those
pairwise slopes divided by \num{1.35} and by the square root of the window's size---the
window's own scatter about its line, expressed as an error on the slope. Over the fitted
range it is \num{0.017}, against a point-to-point variation of \num{0.052}: the curve
moves more between neighbouring lags than the scatter within a window accounts for, which
is the quantitative statement that what the panel shows is structure and not noise. The
half-width is a stated compromise rather than a tuned one: narrower and the amplification
returns, wider and a genuine bend is flattened into the straight line the panel exists to
test for; \num{0.15} decades resolves a crossover of half a decade. It is not a smoother
in disguise. The residual wiggle of this archive is correlated across some four grid
points, so it survives the window---which is the right outcome for something that is not
noise, and the reason the window is reported with the figure rather than chosen to make
the curve look better.
\end{revblock}

\subsection{\rev{The audit pass}}
\begin{revblock}
The ensemble curve turned out not to be a power law (Sec.~\ref{sec:obs-global}), and
none of the questions that decides---is the shape the coverage rule, the survivorship, a
shared heading, or the motion---can be asked of an averaged curve after the fact. Each is
a different reduction of the same per-flight quantity, and that quantity is gone by the
time the average exists.

\path{audit_msd.py} therefore makes a second streaming pass and writes what the first one
discards: one row per flight and one column per lag, holding each flight's $E$ and $N$ at
that lag under the estimator's own coverage rule, \path{NaN} where it does not cover it,
plus a per-flight summary (duration, path, native step, mean-velocity components, largest
in-segment step speed, coordinate extent). The coverage rule is \emph{copied} into that
script rather than imported, so that a change to the estimator surfaces as a disagreement
between the figure and the audit instead of being tracked silently. The arrays run to a
few hundred megabytes per discipline and are analysis products, not thesis ones: they live
outside the repository, under \path{AUDIT_DIR}.

\path{audit_msd_report.py} reduces them, together with the curve table
\path{msd_curve.csv} and the segment and flight tables, into
\texttt{\textbackslash StatAudit*}. Everything the audit asserts in the body reaches it
through those macros. What it computes: the rms residual each estimator leaves against
the power law fitted to it, in dex and as a percentage; the range of the local slope over
the fitted decades; the mean vector's share of the mean square,
$|\langle\mathbf{r}\rangle|^2/\langle|\mathbf{r}|^2\rangle$; the local slope recomputed on
sub-ensembles of a single native cadence, where the population cannot change, and on
fixed-duration cohorts, where it cannot be selected; the first-crossing time of
\SI{5}{\kilo\metre} from take-off, per flight; the per-component variance ratio and the
Kolmogorov--Smirnov distance between the $E$ and $N$ marginals; and the pre-processing
losses of Table~\ref{tab:pipecensus} decomposed by the rule that caused them.

Two of its numbers are worth naming here because they are checks on the pipeline rather
than on the estimator. The retained tables carry
\num{\StatAuditParaNonFiniteFixes} non-finite coordinates and
\num{\StatAuditParaNonMonotoneFlights} non-monotone clocks, which is what
Sec.~\ref{impl:verify} asserts and this pass confirms independently. And the flights whose
first retained segment did not survive resampling start their record at $t>0$ rather than
at the origin---\SI{\StatAuditParaLateStartPct}{\percent} of paraglider and
\SI{\StatAuditHangLateStartPct}{\percent} of hang-glider flights. Their position is still
measured from their own take-off, since the ENU frame is anchored there and a dropped
segment does not move it; they simply do not answer the lags before their record starts.
For the rest, the position at $t=0$ is the origin to within
\SI{\StatAuditParaOriginMedianM}{\metre} in the median, the smoothing's displacement of
the first sample.
\end{revblock}

\subsection{\rev{The shape observables, and what calibrates them}}
\begin{revblock}
\paragraph{\rev{The non-Gaussian parameter.}} \rev{Computed from the moment spectrum
already accumulated, as
$\alpha_2=\langle|\Delta\mathbf{r}|^4\rangle/2\langle|\Delta\mathbf{r}|^2\rangle^2-1$ at
each lag, on order-1 increments. Two calibrations make the archive's value readable and
both are pinned by \path{tests/analysis/observables/test_moments.py}, at the sample size the
test uses: an exact fractional Brownian motion returns \num{0.02} in the median over a
ten-point grid spanning the quoted range, and a Lévy walk of tail index \num{1.5} returns
\num{0.55}. The Gaussian \num{0.02} is estimation bias and not shape --- a fourth moment
over $10^4$ increments is not yet its own expectation --- and sixteen times the data
collapses it to within a hundredth of zero, sign included. The same test pins that too,
since it is what makes the archive's own value a departure rather than an offset. It is quoted over
\SIrange{60}{2000}{\second} rather than over the whole grid; above that the declared task
governs the trajectory and the fourth moment rests on the fewest flights, and the curve
climbs to \StatShapeParaNonGaussBeyond\ there.}

\paragraph{\rev{The tail of the velocity autocorrelation.}} \rev{Green--Kubo is used in
its scaling form, $C(\tau)\sim\tau^{-\gamma}$ with $\gamma<1$ giving
$\langle|\Delta\mathbf{r}|^2\rangle\sim\Delta^{2-\gamma}$, and not as the integral. Two
reasons, both properties of the archive: the integral runs from zero and $C$ is estimable
only above the smoothing scale of the slowest logger, so its first decade is missing; and
$C$ is estimated per flight with that flight's own mean velocity removed, which pulls every
lag down by about the record-mean of $C$ and so steepens the tail. Both push $\gamma$ up,
so $2-\gamma$ is a floor and \path{vacf_tail_exponent} documents it as one. The
one-sidedness is itself tested, on an Ornstein--Uhlenbeck velocity whose true correlation
is known: the estimate with the mean removed is always the steeper.}

\paragraph{\rev{The filter kernel, against a closed form.}} \rev{For any process with
stationary increments and $S(\Delta)=\sigma^2\Delta^{2H}$, the two orders are related
exactly by $V_2=4S(\Delta)-S(2\Delta)=S(\Delta)\,(4-2^{2H})$. That identity pins the
kernel in a way a second convolution cannot: the order-2 kernel is symmetric, so an
independent \path{np.convolve} route performs the same three products in the same order and
agrees to the last bit whether or not the kernel is right --- which is what a
recomputation of the paraglider exponent by that route in fact found. The estimator also
has to be unbiased where the archive sits, at $H$ near \num{1}, which is the boundary of
the family: on exact fractional Brownian motion of $H=0.99$ it returns \num{1.98}.}
\end{revblock}

\paragraph{\rev{When a knee is declared.}} \rev{A bilinear fit to $q\nu(q)$ has two more
parameters than a straight line through the origin, always reduces the residual, and on a
grid of ten moment orders buys the reduction cheaply --- fitted to a fractional Brownian
motion, which has no knee, BIC prefers the bilinear form while the straight line already
fits to an rms of \num{0.003}. So BIC alone is not the rule. \path{bilinear_fit} declares a
knee only when three conditions hold together: BIC prefers the bilinear form, \emph{and} the
straight line leaves a residual above \path{min_departure} = \num{0.02}, \emph{and} the
knee sits inside the $q$ grid rather than on its edge. On the synthetic pair those three
give no knee for fractional Brownian motion and a knee at \num{1.57} with a right-hand slope
of \num{0.96} for a Lévy walk of index \num{1.5}. The archive's departure from a line is
\StatShapeParaLinearDeparture, an order of magnitude below the threshold.}

\paragraph{\rev{Why the drift is filtered and never subtracted.}} \rev{\path{centring_bias}
measures the alternative rather than arguing against it: fractional Brownian motion of known
exponent is generated, the mean-squared displacement is computed raw and after subtracting
each realisation's own net velocity, and the two are compared lag by lag. Subtracting a
velocity estimated from the record's endpoint turns the series into a bridge whose increments
must sum to zero, so the centred curve is pulled down at every lag and hardest at the long
ones --- \SI{2}{\percent} at a lag of \SI{0.05}{\percent} of the record and
\SI{76}{\percent} at half of it, with no range in between where the distortion is
negligible. On $\alpha=1.50$ the operation returns \num{1.35}. The filtered variation of
\path{soaring.analysis.observables.variations} estimates no drift at all, so none can be
estimated badly, and the order scan is what reads off how much drift there was.}

\paragraph{\rev{The variation pass, and why stratifying is free.}} \rev{
\path{measure_variations.py} traverses the fix table once and keeps, per flight and per
filter order, one curve over the lag grid together with the flight's cadence, wing class,
season and declared task. Every stratification of Sec.~\ref{sec:transport-nulls} is then a
row selection on that table rather than another pass over $1.4\times10^{9}$ rows, which is
what makes four cuts affordable where one would otherwise be a day's work.}

\paragraph{\rev{The persistence runs, and the bias that decomposition avoids.}}
\rev{\path{persistence_runs} walks the record greedily from its start, extending a run while
the sinuosity --- arc length over chord --- stays under the threshold and beginning the next
run where the last one ended, with a floor of five samples. The stopping rule is
\emph{first violation}: the condition is not monotone in the endpoint, since a stretch can
come back inside the threshold after leaving it, so the rule is stated rather than implied.
No sample belongs to two runs, and that is the point. Sampling the run length at every
instant instead returns each run in proportion to its own length --- the length-biased
distribution, whose tail index is one below the truth. \path{inspection_biased_runs}
implements the wrong way on purpose so the size of the error can be measured: on synthetic
runs of known index \num{2.0} the two estimators return \num{2.0} and \num{1.1}. At the
values this archive shows that is not a bias to correct afterwards but a different answer.
The first implementation of the biased estimator measured the residual life instead of the
containing run, which is a third quantity again; a test caught it.}

\paragraph{\rev{First-passage times, and why they are not reported.}} \rev{The audit arrays
keep each flight's position at every lag, so the first-passage time $T(L)$ --- the first time
a flight is $L$ from take-off --- costs no traversal, and
$\langle T(L)\rangle\sim L^{1/H}$ would give the exponent in the space domain, independently
of any moment fitted over lags. It was computed and is not reported, for a reason worth
recording so it is not attempted again.
Read naively it returns $\alpha=2.46$, which no process with stationary increments can give.
The cause is right censoring: at \SI{34}{\kilo\metre} only a quarter of paraglider flights
arrive and the missing three quarters are the slow ones. Treating non-arrival as censored and
taking a Kaplan--Meier median instead returns \num{1.95} for paragliders and \num{2.03} for
hang gliders, both inside the paraglider interval --- but refitting over sub-ranges of radius
moves the paraglider value between \num{1.51} and \num{2.03}, a spread five times the
uncertainty it was meant to check.
The reason is the one Sec.~\ref{sec:obs-global} gives for the ensemble MSD. $T(L)$ is measured
from take-off, so it shares that estimator's origin: at small $L$ the wing is still climbing
out inside its launch area, and at large $L$ the censored median extrapolates past the
duration of most flights. It is the same crossover seen along the space axis instead of the
time axis, and it is not an independent check on anything.}

\subsection{\rev{Known limitations, with their size}}
\begin{revblock}
Four, each bounded rather than repaired.

\emph{Edge samples enter.} The smoothing evaluates its polynomial off-centre at the ends
of a segment and the table flags those samples \path{edge} (Sec.~\ref{sec:savgol}).
\path{verify_dataset.py} qualifies its invariants by that flag---which is why its speed
invariant holds---and the MSD estimators do not, and read every fix. Those samples are
under \SI{1}{\percent} of the table and account for
\num{\StatVerifyParaEdgeStepsOver} of the \num{\StatVerifyParaEdgeStepsOver}${}+{}$%
\num{\StatVerifyParaMeasuredStepsOver} in-segment steps above the fix-level speed bound,
\SI{\StatVerifyParaEdgeStepsPct}{\percent} of all steps. The worst step between two
\emph{measured} samples is \SI{\StatVerifyParaWorstMeasuredMs}{\metre\per\second} against
a bound of \SI{45}{\metre\per\second}: an excess of under a tenth, which the verifier
reports and does not fail on, its threshold for a data defect being twice the bound.
Excluding the flagged samples moves the ensemble MSD by up to \SI{13}{\percent} below
\SI{10}{\second} and by under \SI{0.2}{\percent} above \SI{60}{\second}, so the fit ranges
are unaffected; excluding them would also cost every flight its shortest lags, since the
first samples of a record are exactly the flagged ones.

\emph{Absolute position carries re-acquisition offsets.} Bounded in
Sec.~\ref{sec:obs-global} at under a per cent of the mean over the fitted range, on
\SI{\StatVerifyParaOffsetPct}{\percent} of paraglider flights.

\emph{The archive holds a few duplicate flights.} The same physical flight uploaded
twice, detected by take-off coordinates, duration and path length agreeing to four
decimal places, one second and ten metres:
\num{\StatAuditParaDupFlights} paraglider flights in
\num{\StatAuditParaDupGroups} groups, \SI{\StatAuditParaDupPct}{\percent} of the archive,
and \num{\StatAuditHangDupFlights} hang-glider flights in \num{\StatAuditHangDupGroups}
groups, \SI{\StatAuditHangDupPct}{\percent}. Ensemble averaging
assumes independent draws, and at this rate the violation is immaterial; it is reported so
that it is not assumed to be zero.

\emph{The last lag of the grid is empty.} \SI{43200}{\second} is the flight-level
duration bound, so no retained flight reaches it and the row is written with a null MSD
and a zero count. It is dropped by every fit and every plot through the finiteness mask,
and appears in \path{msd_curve.csv} for grid completeness.
\end{revblock}

\paragraph{\rev{Reading the table.}} \rev{Through
\path{soaring.analysis.derived.stream_flights}, never by iterating Parquet row groups: a
row group is a unit of storage and cuts across flights, so the direct reading counted
\SI{0.7}{\percent} of flights twice and truncated the longest segments, which are
precisely the ones the long-lag end of the time-averaged curve rests on.}
